\subsection{Unprotected targets and a four-unit residual margin}

Continue with a finite nonempty, strongly connected, arc-minimal
all-deficient oriented graph. Here arc-minimal means minimal under
deletion of arbitrary nonempty sets of arcs on the fixed vertex set,
not merely under single-arc deletion.
Write $\mathcal F(x)$ for the vertices $y\in U(x)$ which are not
protected, and put $t(x)=|\mathcal F(x)|$. Protection means
\[
 y\to x,\qquad \Nin(y)\subseteq\Nin(x)\setminus\{y\}.
\]
All neighbourhoods, convenient orientations, and protected pairs below
refer to the original graph.

\begin{lemma}[The unprotected-target envelope]
\label{lem:unprotected-envelope}
Fix a vertex $x$ with $P=P(x)\ne\varnothing$, and put
\[
 \begin{gathered}
 Q=\Nback(x),\quad W=\Nin(x),\quad C=\bd P,\quad H=Q\setminus C,\\
 X=\bd^2P,\quad j=|H|,\quad s=|\Nout(x)\cap H|.
 \end{gathered}
\]
Then
\[
 j=s=0\Longrightarrow t(x)\le\eta(x),\qquad
 j=s>0\Longrightarrow t(x)\le\eta(x)-g(x).
\]
\end{lemma}
\begin{proof}
When $j=s$, every vertex of $H$ is an original outneighbour of $x$.
Delete all arcs $x\to H$, and call the resulting graph $\widehat D$;
if $H=\varnothing$, set $\widehat D=D$. Its retained first neighbours
lie in $P\cup C$. Also
$\mathcal M(x)\subseteq P\cup C$, whereas $X$ is disjoint from $P\cup C$.
The target-deletion envelope, including deleted first neighbours that
may become exact second neighbours, gives
\[
 \Ntwo_{\widehat D}(x)\subseteq
 E:=\mathcal M(x)\,\dot\cup\,X,
 \qquad
 |E|\le\mu(x)+|Q|-s-1
       =\dtwo_D(x)+\eta(x)-s.
\]
For completeness, if a retained intermediate is in $P$, its outgoing
arcs stay in $P\cup C$; if it is in $C$, an outgoing head outside that
set is in $X$. Inside $P\cup C$, a head which is not a retained first
neighbour is an original nonneighbour of $x$. A deleted head in $H$
which becomes second belongs to $X$ and is counted only once.
The size bound uses $|X|\le|C|-1=|Q|-s-1$.

We claim
\[
 \mathcal F(x)\subseteq E\setminus\Ntwo_{\widehat D}(x).
\]
For a missing far target this follows from its membership in
$\mathcal M(x)$. Otherwise a far target $y$ lies in $W$ and points to
$x$. Being unprotected, it has a predecessor $z\notin W$.
The vertex $z$ cannot be $x$; it cannot be in $P$, because no arc from
$P$ points into $W$; and it cannot be in $H\subseteq\Nout(x)$, because
that would give $x\to z\to y$. Thus $z\in C$ and $y\in X$.
Finally, deleting arcs cannot create a path of length at most two to
an originally unreachable target, so $y\notin\Ntwo_{\widehat D}(x)$.

If $s=0$, subtracting the original second neighbourhood from $E$
gives $t(x)\le\eta(x)$. If $s>0$, every other vertex has unchanged
first neighbourhood and cannot gain second neighbours. Arc-minimality
therefore makes $x$ a Seymour vertex in $\widehat D$, so
$\dtwo_{\widehat D}(x)\ge\dout_D(x)-s$. Consequently
\[
 t(x)\le|E|-\dtwo_{\widehat D}(x)
 \le\dtwo_D(x)+\eta(x)-s-(\dout_D(x)-s)
 =\eta(x)-g(x).
\]
No minimality or counterexample property is assumed for $\widehat D$.
\end{proof}

\begin{corollary}[One unit of residual]\label{cor:eta-one-budget}
If $\eta(x)=1$, then $t(x)\le1$.
\end{corollary}
\begin{proof}
If $P(x)=\varnothing$, the identity
$1=2\mu(x)+g(x)-1$ and $g(x)\in\{1,2\}$ force
$\mu(x)=0$ and $g(x)=2$. For a whole vertex every far target is
protected: it points to $x$, and each of its predecessors must also
point to $x$, because the opposite orientation would give a two-step
path from $x$ to that target. Thus $t(x)=0$.

If $P(x)\ne\varnothing$, the existing local bounds are
\[
 s=0\Longrightarrow\eta(x)\ge2j,\qquad
 s>0\Longrightarrow\eta(x)\ge g(x)+2(j-s).
\]
For residual one, the first case forces $j=s=0$. The second forces
$g(x)=1$ and $j=s>0$. Apply
Lemma~\ref{lem:unprotected-envelope}; in the second case it even gives
$t(x)=0$.
\end{proof}

\begin{lemma}[A two-unit dichotomy]\label{lem:eta-two-dichotomy}
Suppose $\eta(x)=2$ and $r$ is a missing partner of $x$ with
$\Nout(r)\cap U(x)\ne\varnothing$. Then either $t(x)\le2$, or
$r\to x$ is a convenient orientation of the missing edge $rx$.
\end{lemma}
\begin{proof}
If $P(x)=\varnothing$, every other vertex reaches $x$ within two steps,
which proves convenience of $r\to x$ immediately.
Otherwise use the notation of Lemma~\ref{lem:unprotected-envelope}.
If $s>0$, the bound
$2\ge g(x)+2(j-s)$ forces $j=s$, and that lemma gives
$t(x)\le2-g(x)\le1$. If $s=j=0$, it gives $t(x)\le2$.
The only remaining case is $s=0,j=1$. Write $H=\{q\}$.
Since $q\in Q$ but $q\notin\Nout(x)$, it is an original missing
partner of $x$.

The retained original first neighbours lie in $P\cup C$, and the
same envelope argument gives
\[
 \Ntwo_D(x)\subseteq
 E_0:=\bigl(\mathcal M(x)\setminus\{q\}\bigr)\,\dot\cup\,X.
\]
The boundary inequality gives $|X|\le|Q|-2$, hence
\[
 |E_0|\le\mu(x)-1+|Q|-2
          =\dtwo_D(x),
\]
where the last equality uses $\eta(x)=2$. Thus the envelope is
saturated:
\[
 \Ntwo_D(x)=E_0.
\]
Every vertex in $P\cup C$ is reachable from $x$ within at most two
steps: it is either an original first neighbour, or a missing partner
other than $q$ and hence an exact second neighbour by saturation.
This assertion is about reachability within two steps, not about
membership of all these vertices in the exact second neighbourhood.
Every arc leaving $P\cup C$ has its head in $X$, also an exact second
neighbour of $x$. A missing partner $r\ne q$ lies in $P\cup C$,
so none of its outneighbours can lie in $U(x)$. The hypothesis on $r$
therefore forces $r=q$.

Finally $q\notin P$ and $q\notin C=\bd P$, so no vertex of $P(x)$
points to $q$. Every predecessor of $q$ is thus outside $P(x)$ and
reaches $x$ within two steps. This is precisely convenience of
$q\to x$.
\end{proof}
